\documentclass[11pt,a4paper]{article}

\usepackage[T1]{fontenc}
\usepackage[utf8]{inputenc}
\usepackage[shorthands=off,turkish]{babel}
\usepackage[a4paper,margin=2.4cm]{geometry}
\usepackage{lmodern}
\usepackage{amsmath}
\usepackage{booktabs}
\usepackage{longtable}
\usepackage{array}
\usepackage{enumitem}
\usepackage{xcolor}
\usepackage{fancyhdr}
\usepackage[hidelinks]{hyperref}

\definecolor{hazir}{HTML}{04704D}
\definecolor{kismi}{HTML}{C07807}
\definecolor{eksik}{HTML}{B91C3C}
\definecolor{kod}{HTML}{2F3735}

\newcommand{\Hazir}{\textcolor{hazir}{\textbf{Hazır}}}
\newcommand{\Kismi}{\textcolor{kismi}{\textbf{Kısmi}}}
\newcommand{\Eksik}{\textcolor{eksik}{\textbf{Eksik}}}
\newcommand{\dosya}[1]{\textcolor{kod}{\texttt{#1}}}

\pagestyle{fancy}
\fancyhf{}
\lhead{\small NLP Mimarisi ve Yöntem Haritası}
\rhead{\small TEKNOFEST 2026 --- 2. Senaryo}
\cfoot{\thepage}

\title{\vspace{-1.5cm}\textbf{NLP Mimarisi ve Yöntem Haritası}\\[0.3cm]
\large Katılım Bankacılığı Kampanya Metinlerinden Bilgi Çıkarımı\\
\normalsize Doğal dil işleme yöntemlerinin şartname maddeleriyle eşlenmesi}
\author{}
\date{\today}

\begin{document}
\maketitle
\thispagestyle{fancy}
\vspace{-1cm}

\section{Amaç ve Kapsam}

Bu belge, geliştirilen çözümde \textbf{doğal dil işlemenin tam olarak nerede ve
hangi yöntemle kullanıldığını} tanımlar. Amaç, şartnamenin ``NLP yöntemlerinin
probleme uygun kullanılması'', ``metin ön işleme adımlarının doğru uygulanması''
ve ``bilgi çıkarımı yaklaşımının başarısı'' değerlendirme maddelerine
izlenebilir bir karşılık vermektir.

\subsection{Kapsam dışı olanlar}

Dürüstlük açısından önce sınır çizilmelidir. Aşağıdaki bileşenler yazılım
mühendisliğidir, doğal dil işleme değildir ve bu belgede NLP olarak
sunulmamaktadır:

\begin{itemize}[leftmargin=1.2cm,itemsep=2pt]
  \item Karşılaştırma motoru ve sıralama mantığı (aritmetik normalizasyon)
  \item Dashboard, arayüz ve görselleştirme katmanı
  \item Veri modeli, kalıcı depo ve API sözleşmesi
  \item On-premise dağıtım mimarisi
\end{itemize}

Doğal dil işleme, yalnızca Bölüm~\ref{sec:hat}'te tanımlanan sekiz adımlık
metin işleme hattında yer alır.

\section{Şartname Maddeleriyle Eşleme}

\begin{longtable}{p{1.1cm} p{3.5cm} p{4.6cm} p{2.9cm} c}
\toprule
\textbf{Madde} & \textbf{İstenen} & \textbf{NLP yöntemi} & \textbf{Modül} & \textbf{Durum} \\
\midrule
\endhead
5.1 & Veri toplama & Gövde metni çıkarma, kalıp (boilerplate) temizliği & --- & \Eksik \\
5.2 & Farklı ve örtük ifade biçimleri & Eşanlamlı kalıp sözlüğü + kural desenleri + dil modeli & \dosya{lexicon.ts} & \Kismi \\
5.3 & Finansal bilgi çıkarımı & Bilgi çıkarımı: değer tanıma, birim çözümleme, aralık ve olumsuzluk kapsamı & \dosya{extract.ts} & \Kismi \\
5.4 & Kampanya türü belirleme & Metin sınıflandırma (TF--IDF + doğrusal model) & --- & \Eksik \\
5.5 & Katılım terminolojisi & Sözlükbirim tabanlı normalizasyon, çekim ekine tolerans & \dosya{lexicon.ts} & \Hazir \\
5.6 & Standart formata dönüştürme & Sayı, para birimi ve tarih normalizasyonu; yazılı sayı ayrıştırma & \dosya{normalize.ts} & \Hazir \\
5.8 & Veri ön işleme & Temizleme, normalizasyon, cümle bölütleme, belirteçleme & \dosya{normalize.ts}, \dosya{segment.ts} & \Hazir \\
--- & Kanıt zorunluluğu & Cümle düzeyinde hizalama (birebir / normalize / örtüşme) & \dosya{align.ts} & \Hazir \\
--- & Chatbot & Niyet sınıflandırma + slot doldurma + sorguya çevirme & --- & \Eksik \\
--- & Model başarısı & Alan bazlı kesinlik / duyarlılık / F1 ölçümü & --- & \Eksik \\
\bottomrule
\end{longtable}

\section{İşleme Hattı}
\label{sec:hat}

Metin, kaynaktan şemaya kadar sekiz adımdan geçer. Adım sırası tesadüfi
değildir: kural tabanlı katman dil modelinden \emph{önce} çalışır, böylece
kolay vakalar deterministik olarak çözülür ve dil modeli yalnızca varyasyonlu
ifadeler için devreye girer.

\begin{enumerate}[leftmargin=1.4cm,itemsep=5pt]
  \item \textbf{Ön işleme} --- \dosya{normalize.ts}\\
        Gövde metni çıkarma, Unicode NFC, tipografik tırnak ve tire
        sadeleştirme, bölünmez boşluk temizliği, Türkçe yerel ayarlı küçültme,
        ASCII katlama.
  \item \textbf{Cümle bölütleme} --- \dosya{segment.ts}\\
        Ondalık ayracı, tarih, binlik ayracı ve kısaltma korumalı bölütleme;
        her cümlenin kaynak metindeki karakter aralığı saklanır.
  \item \textbf{Terminoloji eşleme} --- \dosya{lexicon.ts}\\
        Konvansiyonel terimlerin katılım karşılıklarına eşlenmesi;
        \texttt{terim\_esleme\_uygulandi} bayrağının belirlenmesi.
  \item \textbf{Kural tabanlı bilgi çıkarımı} --- \dosya{extract.ts}\\
        Oran ve periyot, vade, tahsis ücreti (olumsuzluk dâhil), tutar aralığı.
  \item \textbf{Metin sınıflandırma} --- \emph{planlanan}\\
        Kampanya türünün sekiz sınıftan birine atanması.
  \item \textbf{Dil modeli} --- yerel, şema kısıtlı üretim\\
        Yalnızca dördüncü adımda boş kalan alanlar için çağrılır.
  \item \textbf{Çapraz doğrulama} --- \dosya{extract.ts}\\
        Kural ve model çıktılarının karşılaştırılması; uyuşmazlıkta güven
        skorunun düşürülmesi.
  \item \textbf{Kanıt hizalama} --- \dosya{align.ts}\\
        Her alanın dayandığı cümlenin kaynak metinde konumlandırılması.
\end{enumerate}

\section{Yöntemlerin Ayrıntısı}

\subsection{Türkçe metin normalizasyonu}

JavaScript'in yerleşik \texttt{toLowerCase()} metodu Türkçe alfabesini
tanımaz. Büyük İ harfi küçültülürken \texttt{i} artı ayrı bir birleşik nokta
karakterine (U+0307) dönüşür; noktasız I ise ASCII \texttt{i} olur:

\begin{center}
\begin{tabular}{lll}
\toprule
\textbf{Girdi} & \textbf{\texttt{toLowerCase()}} & \textbf{Sonuç} \\
\midrule
FAİZ ORANI & \texttt{fai}\,\textcolor{eksik}{\texttt{[U+0307]}}\,\texttt{z oranı} & Terim yakalanamaz \\
DOSYA MASRAFI & \texttt{dosya masraf}\textcolor{eksik}{\texttt{i}} & Terim yakalanamaz \\
\bottomrule
\end{tabular}
\end{center}

Bu nedenle tüm küçültme işlemleri \texttt{tr-TR} yereliyle yapılır ve
eşleştirme, aksan ile şapkayı katlayan ayrı bir anahtar üzerinden yürütülür.
Böylece ``kâr payı'' ile ``kar payi'' aynı terime çözümlenir.

Sayısal normalizasyon üç biçimi tek değere indirger: Türkçe ondalık
(\texttt{2,05}), İngilizce ondalık (\texttt{2.05}) ve binlik ayraçlı
(\texttt{1.000.000}). Ayrıca yazıyla yazılmış Türkçe sayılar çözümlenir;
``beş yüz Türk Lirası dosya masrafı'' ifadesi \texttt{500} değerini üretir.

\subsection{Cümle bölütleme}

Naif nokta bölme finans metinlerinde başarısız olur. Bölütleyici şu kalıpları
cümle sonu saymaz: ondalık ayracı (\texttt{\%2,05}), tarih
(\texttt{31.12.2026}), binlik ayracı (\texttt{1.000.000}), Türkçe kısaltmalar
(\texttt{vb.}, \texttt{Sn.}, \texttt{md.}) ve baş harf kullanımı
(\texttt{A. Bankası}). Her cümle için \((\text{başlangıç}, \text{bitiş})\)
karakter aralığı üretilir; kanıt hizalama bu aralıklara dayanır.

\subsection{Sözlükbirim eşleme ve morfolojik tolerans}

Türkçe eklemeli bir dildir; sabit kelime araması yetersizdir. Sözlük
girdileri kök hâlinde tutulur ve eşleştirme, kökün ardından en fazla altı
harflik çekim ekine izin veren bir desenle yapılır. Böylece \emph{faiz},
\emph{faizi}, \emph{faizde}, \emph{faizli} tek bir girdiyle yakalanır.

\subsection{Olumsuzluk tespiti}

``Tahsis ücreti alınmaz'' ifadesi şemada \texttt{0} değerine karşılık gelir;
pozitif bir sayı aranarak bulunamaz. Türkçede olumsuzluk çekim ekiyle
taşınır:

\begin{center}
\begin{tabular}{lll}
\toprule
\textbf{Ek} & \textbf{Örnek} & \textbf{Değer} \\
\midrule
-maz / -mez & alınmaz, edilmez & 0 \\
-mıyor / -muyor & alınmıyor & 0 \\
-mama- / -meme- & alınmamaktadır & 0 \\
\midrule
-makta / -mekte & alınmaktadır & pozitif \\
\bottomrule
\end{tabular}
\end{center}

Son satır kritiktir: ``alınmaktadır'' \emph{olumludur}. Naif bir ``alınm''
araması bu iki durumu ayırt edemez ve ücreti hatalı biçimde sıfırlar.

\subsection{Kural tabanlı bilgi çıkarımı}

Birinci katman deterministik, milisaniyelik ve tekrarlanabilirdir. Dört alanı
bağlam denetimiyle çıkarır: oran ifadesinin gerçekten kâr payına ait olup
olmadığı, vade ifadesinin taksit sayısıyla karışıp karışmadığı, ücret
cümlesinin tutar cümlesinden ayrılması gibi denetimler uygulanır.

Şartname kuralı doğrudan kodlanmıştır: metinde ``aylık'' veya ``yıllık''
geçmiyorsa periyot \texttt{belirsiz} atanır ve o alanın güven skoru en fazla
0,5 olur.

\subsection{Kanıt hizalama}

Modelin döndürdüğü alıntı, kaynak metinle birebir aynı olmayabilir. Hizalayıcı
üç aşamalıdır: birebir arama, normalize edilmiş metin üzerinde arama, ve
cümle düzeyinde belirteç örtüşmesi. Üçüncü aşamada Jaccard benzerliği
kullanılır:

\[
J(A,B) \;=\; \frac{|A \cap B|}{|A \cup B|}
\]

Burada \(A\) alıntının, \(B\) aday cümlenin belirteç kümesidir. Eşik değeri
\(J \geq 0{,}45\) olarak belirlenmiştir. Çakışan aralıklarda yüksek skorlu
hizalama kazanır.

\subsection{Metin sınıflandırma (planlanan)}

Kampanya türü sekiz sınıfa ayrılacaktır. Yöntem, belirteç ağırlıklandırması
ve doğrusal bir sınıflandırıcıdır:

\[
\mathrm{tfidf}(t,d) \;=\; \mathrm{tf}(t,d)\cdot\log\frac{N}{\mathrm{df}(t)}
\]

Bu tercih bilinçlidir: yerelde saniyeler içinde eğitilir, GPU gerektirmez,
kararı açıklanabilir (hangi terimin hangi sınıfa katkı verdiği görülebilir)
ve on-premise kısıtıyla uyumludur.

\subsection{Çapraz doğrulama}

Kural katmanı ile dil modeli bağımsız iki kaynaktır. Değerler \%1 tolerans
içinde uyuşuyorsa güven korunur; uyuşmuyorsa güven skoru yarıya indirilir ve
kayıt manuel doğrulama kuyruğuna alınır. Bu, tek bir modelin kendi beyan
ettiği güven skoruna bağımlı kalmayı önler.

\subsection{Chatbot için dil işleme (planlanan)}

Chatbot, serbest üretim yerine yapılandırılmış veriyi sorgular. Gerekli dil
işleme adımları:

\begin{itemize}[leftmargin=1.2cm,itemsep=2pt]
  \item \textbf{Niyet sınıflandırma:} tek ürün sorgusu, karşılaştırma, listeleme, tanım sorusu
  \item \textbf{Slot doldurma:} banka adı, ürün türü, kriter, sayısal kısıt
  \item \textbf{Sorguya çevirme:} slotlardan veritabanı sorgusu üretimi
  \item \textbf{Cevap üretimi:} şablon tabanlı, kaynak cümle ve güven skoru ekli
\end{itemize}

\section{Değerlendirme Kriterleriyle İlişki}

\begin{center}
\begin{tabular}{p{6.6cm} p{7.2cm}}
\toprule
\textbf{Kriter} & \textbf{Karşılayan bileşen} \\
\midrule
NLP yöntemlerinin probleme uygun kullanılması & Katmanlı hat: kural $\rightarrow$ model $\rightarrow$ çapraz doğrulama \\
Metin ön işleme adımlarının doğru uygulanması & \dosya{normalize.ts}, \dosya{segment.ts} \\
Bilgi çıkarımı yaklaşımının başarısı & \dosya{extract.ts} + ölçüm seti \\
Farklı ifade biçimlerini doğru yorumlayabilmesi & Morfolojik tolerans, eşanlamlı kalıp sözlüğü, dil modeli \\
Eksik/farklı yazılmış bilgide doğru sonuç & Olumsuzluk tespiti, yazılı sayı çözümleme, ASCII katlama \\
\bottomrule
\end{tabular}
\end{center}

\section{Ölçüm Planı}

Her alan için ayrı ayrı hesaplanacaktır:

\[
P = \frac{TP}{TP+FP}, \qquad
R = \frac{TP}{TP+FN}, \qquad
F_1 = \frac{2PR}{P+R}
\]

Bir alan, altın değerle tam eşleştiğinde doğru sayılır; sayısal alanlarda
\%1 tolerans uygulanır. Ölçüm üç yapılandırmada tekrarlanır: yalnızca kural
katmanı, yalnızca dil modeli, ve katmanlı hibrit. Bu karşılaştırma, hibrit
mimarinin katkısını sayısal olarak gösterir.

\section{Eksikler ve İş Planı}

\begin{center}
\begin{tabular}{p{7.2cm} p{2.6cm} p{3.4cm}}
\toprule
\textbf{İş} & \textbf{Süre} & \textbf{Karşıladığı madde} \\
\midrule
Kural motorunu birinci katmana taşıma & Yarım gün & 5.3, teknik kriter \\
Örtük ifade kalıp sözlüğü & Birkaç saat & 5.2 \\
Kampanya türü sınıflandırıcı & Yarım gün & 5.4 \\
Ölçüm seti ve F1 raporu & Yarım gün & Model başarısı \\
Ön işleme hattının toplama katmanına bağlanması & Yarım gün & 5.1, 5.8 \\
Chatbot niyet ve slot katmanı & 1,5 gün & Senaryo kapsamı \\
\bottomrule
\end{tabular}
\end{center}

\vspace{0.4cm}
\noindent
Toplam yaklaşık dört iş günü. Sıralamada ilk madde önceliklidir: kural
katmanının öne alınması hem gecikmeyi düşürür hem de sonraki ölçümler için
temel oluşturur.

\end{document}
